\documentclass[11pt]{article}
\usepackage[a4paper,margin=25mm]{geometry}
\usepackage{amsmath,amssymb,amsthm,hyperref}
\newtheorem{theorem}{Theorem}
\emergencystretch=2em
\title{A nonexpansive semigroup's Cesaro limit need not be the nearest fixed point\\Disproof of conjecture 00000008856}
\author{ziangni-sys}
\date{4 October 2026}
\begin{document}
\maketitle
\paragraph{Claim and indexing convention.}
Both source versions assert that Cesaro average orbits of nonexpansive semigroups converge weakly to the metric projection of the starting point onto the common fixed-point set. We use a \emph{discrete} semigroup indexed by the additive monoid $\mathbb N$, generated by a single nonexpansive map on the real Hilbert space. The source describes nonexpansive map families and does not restrict to continuous-time semigroups or impose continuity in a real time parameter. Our example concerns its unrestricted statement, and makes no assertion about a narrower continuous-time theorem.

\paragraph{An actual nonexpansive semigroup.}
On $\mathbb R$ with its usual Euclidean metric set
\[
 T_0(x)=x,\qquad T_n(x)=|x|\quad(n\ge1).
\]
The map $T(x)=|x|$ is idempotent, so $T_n=T^n$ and
$T_{m+n}=T_m\circ T_n$ for all natural $m,n$. Every map is nonexpansive: the identity is isometric and
\[
 \bigl||x|-|y|\bigr|\le |x-y|.
\]
The common fixed-point set is exactly
\[
 F=\{x:T_nx=x\text{ for every }n\}=[0,\infty).
\]
It is nonempty, closed and convex. This avoids any failure caused by an empty fixed-point set or an unbounded orbit.

\begin{theorem}
For $x_0=-1$, the Cesaro averages converge strongly and weakly to $1$, whereas the metric projection of $x_0$ onto $F$ is $0$.
\end{theorem}
\begin{proof}
The orbit is $-1,1,1,\ldots$ and is bounded in absolute value by one. Averaging the first $N\ge1$ terms gives
\[
 a_N=\frac1N\sum_{k=0}^{N-1}T_k(-1)=\frac{-1+(N-1)}N
     =1-\frac2N\longrightarrow1.
\]
Strong convergence implies weak convergence. Any weak limit must also be $1$: applying the defining weak-convergence condition to the identity continuous linear functional on $\mathbb R$ yields ordinary convergence, whose limit is unique.

For $z\in F$, the distance from $-1$ is $|-1-z|=1+z$, minimized uniquely at $z=0$. Thus the actual metric projection is zero. The weak limit one has distance two from the initial point and cannot be its nearest fixed point. No point can simultaneously be a weak limit of this average orbit and a nearest fixed point.
\end{proof}

\newpage
\paragraph{Scope of the disproof.}
This example refutes the projection identification in the first conjectured clause, while the averages themselves converge. Neither absence of fixed points nor an unbounded trajectory is involved. The later claims concerning angular density, linear rates and strong convergence criteria are separate conjuncts and do not need separate disproofs. The fixed set is a half-line, rather than a linear subspace, and the generating map is nonlinear. A projection-identification result for linear contraction semigroups therefore would not contradict this example.

Discrete semigroups are a standard class of nonexpansive map families and Cesaro averages are the finite averages displayed above. The family is indexed by $\mathbb N$; it is not presented as a strongly continuous family indexed by all real times. This distinction is part of the stated scope.

\paragraph{Faithful Lean correspondence.}
The formalization uses the actual real Hilbert space with its actual distance and absolute value. Theorems \path{semigroup_identity}, \path{semigroup_law} and \path{actual_iterates} verify the identity element, all additive-index compositions and the relation to iterates of the generating map. \path{nonexpansive} proves the metric inequality for every map and every pair of real inputs.

The actual common fixed set quantifies over all indices, and is proved equal to \path{Set.Ici 0}. Its closedness, convexity and nonemptiness are proved. The orbit boundedness theorem covers every index.

The function \path{average} is the actual finite sum of the first $n+1$ iterates divided by their count. \path{actual_sum} proves the numerator for every $n$ by induction; \path{average_formula} derives $1-2/(n+1)$. The theorem \path{average_strong_limit} is actual topological convergence at infinity using the standard reciprocal-natural-number limit.

Weak convergence is defined by quantifying over every actual continuous linear functional. \path{average_weak_limit} derives it from strong convergence; \path{weak_limit_unique} uses the actual identity continuous linear map. This is the full standard functional criterion, not a claim that a separate weak-topology API was used.

The predicate \path{Nearest} requires membership and a distance inequality against every fixed point. \path{actual_metric_projection} proves that zero has this property, and \path{metric_projection_unique} shows every point with it equals zero. The final \path{no_nearest_weak_limit} negates the existence of a weak limit satisfying the nearest-point property. \path{counterexample} packages the semigroup, nonexpansiveness, nonempty fixed set, bounded orbit and obstruction together. None of the orbit, average, fixed set, norm or projection values is assumed as a supplied certificate.

\paragraph{Reproduction.}
The public project pins Lean 4.19.0 and Mathlib commit:\\
\path{c44e0c8ee63ca166450922a373c7409c5d26b00b}.
Run \texttt{lake build} from \path{lean/}. Principal theorem axiom audits are printed by the source. No admitted proof, custom axiom or computation-trust shortcut is used.
\end{document}
